% sections/02-background.tex

\section{Background}
\label{sec:background}

\subsection{The Kulldorff Spatial Scan Statistic}
\label{sec:background:kulldorff}

Let $S = \{s_1, \ldots, s_K\}$ be a set of $\Ktot$ reporting nodes (e.g., provinces or
sampling sites), each associated with an observed case count $c_i$ and a testing
volume $n_i$.
Let $\Ctot = \sum_i c_i$ be the national total of observed cases and $\Ntot = \sum_i n_i$
the national total testing volume.

A \emph{candidate window} $w$ is a pair (circle, time period) parameterized by a center
node, a spatial radius, and a temporal duration.
For each window $w$, let $\kw$ be the number of nodes inside $w$,
$\nw = \sum_{i \in w} n_i$ the testing volume inside $w$,
and $\cobs = \sum_{i \in w} c_i$ the observed cases inside $w$.
The expected case count under the null hypothesis of uniform distribution is:
\begin{equation}
    \muexp = \Ctot \cdot \frac{\nw}{\Ntot}
    \label{eq:mu}
\end{equation}

The \emph{log-likelihood ratio} (LLR) for window $w$ under the Poisson model
is~\cite{kulldorff1997}:
\begin{equation}
    \LLR(w) =
    \begin{cases}
        \cobs \ln\!\left(\dfrac{\cobs}{\muexp}\right)
        + (\Ctot - \cobs) \ln\!\left(\dfrac{\Ctot - \cobs}{\Ctot - \muexp}\right)
        & \text{if } \cobs > \muexp \\[6pt]
        0 & \text{otherwise}
    \end{cases}
    \label{eq:llr}
\end{equation}

The \emph{relative risk} inside $w$ is:
\begin{equation}
    \RR(w) = \frac{\cobs / \nw}{(\Ctot - \cobs) / (\Ntot - \nw)}
    \label{eq:rr}
\end{equation}

Note that $\kw$ does not appear in \cref{eq:llr}: the detection question requires
only aggregate counts.
The privacy question, however, requires knowing how many independent actors contributed
to those aggregates, which is exactly $\kw$.

The primary cluster is $\wstar = \arg\max_{w \in \Wcal} \LLR(w)$.
Statistical significance is assessed by Monte Carlo simulation: $M = 999$ datasets
are drawn under the null hypothesis, and the p-value is
$p = (\#\{m : \LLR_m \geq \LLR(\wstar)\} + 1) / (M + 1)$~\cite{kulldorff1997}.

Algorithm~\ref{alg:kulldorff} formalizes the complete procedure.

% algorithms/alg1-kulldorff.tex
% Algorithm 1: Non-private Kulldorff scan statistic (established, Kulldorff 1997)
% Placed in background section. The private variant is alg2-private.tex.

\begin{algorithm}[H]
\caption{Kulldorff Prospective Space-Time Scan Statistic~\cite{kulldorff1997,neill2004}}
\label{alg:kulldorff}
\begin{algorithmic}[1]

\Statex \textbf{Input:}
  Spatial units $S = \{s_1, \ldots, s_\Ktot\}$;
  counts $c_i$ and testing volumes $n_i$ for each $s_i$;
  national totals $\Ctot$, $\Ntot$;
  max spatial window $\alpha$ (\% of $\Ntot$);
  max temporal window $T_{\max}$;
  simulations $M$
\Statex \textbf{Output:}
  Primary cluster $\wstar$ with $\LLR(\wstar)$, $\RR(\wstar)$, $p$-value

\Statex
\Statex \textit{// Step 1: Score every candidate window}
\State $\Wcal \gets \emptyset$
\For{each center $s_i \in S$}
    \For{each radius $r$ by nearest-neighbor, while $\nw \leq \alpha \cdot \Ntot$}
        \For{$t = 1, \ldots, T_{\max}$}
            \State $w \gets (s_i,\; r,\; t)$
            \State $\cobs \gets \sum_{s_j \in w} c_j$
                \Comment{observed cases inside $w$}
            \State $\nw \gets \sum_{s_j \in w} n_j$
                \Comment{testing volume inside $w$}
            \State $\muexp \gets \Ctot \cdot \nw / \Ntot$
                \Comment{expected cases under null}
            \If{$\cobs > \muexp$}
                \State $\LLR(w) \gets
                    \cobs \ln\!\left(\dfrac{\cobs}{\muexp}\right)
                    + (\Ctot - \cobs)\ln\!\left(\dfrac{\Ctot - \cobs}{\Ctot - \muexp}\right)$
                \State $\RR(w) \gets
                    \dfrac{\cobs / \nw}{(\Ctot - \cobs) / (\Ntot - \nw)}$
            \Else
                \State $\LLR(w) \gets 0$
            \EndIf
            \State $\Wcal \gets \Wcal \cup \{w\}$
        \EndFor
    \EndFor
\EndFor

\Statex
\Statex \textit{// Step 2: Select primary cluster (deterministic argmax)}
\State $\wstar \gets \arg\max_{w \in \Wcal}\; \LLR(w)$

\Statex
\Statex \textit{// Step 3: Significance testing via Monte Carlo}
\State $\textit{count} \gets 0$
\For{$m = 1$ to $M$}
    \Comment{$M = 999$ by convention}
    \State Redistribute $\Ctot$ cases across $S$ proportional to $n_i$
    \State $\LLR_m \gets \max_{w \in \Wcal}\; \LLR(w)$ on randomized data
    \If{$\LLR_m \geq \LLR(\wstar)$}
        \State $\textit{count} \gets \textit{count} + 1$
    \EndIf
\EndFor
\State $p \gets (\textit{count} + 1) / (M + 1)$

\State \Return $\wstar,\; \LLR(\wstar),\; \RR(\wstar),\; p$

\end{algorithmic}
\end{algorithm}


% ── Data structure example ────────────────────────────────────────────────────

\paragraph{Running example.}
Table~\ref{tab:data} shows a synthetic dataset on Vietnamese province structure
used throughout this paper.
The sensitivity column ($\sens$) does \emph{not} appear here; it is the quantity
derived in Theorem~\ref{thm:sensitivity}.

% tables/tab1-data-structure.tex
% Table 1: Synthetic dataset showing the data structure.
% NOTE: Sensitivity column (Delta(k_w)) does NOT appear here —
%       it is the quantity derived in Theorem 1 (contribution section).
% Values are synthetic. mu = C * n_w / N must verify correctly for every row.
% Using C = 100, N = 10,000.

\begin{table}[ht]
\centering
\caption{Synthetic dataset illustrating the scan statistic data structure
         (C = 100 total cases, N = 10{,}000 total samples).
         Window W\textsubscript{02} is included in the table but is not shown
         on the map in Figure~\placeholder{X}; the table is illustrative.
         The sensitivity column $\sens$ is absent here; it is the quantity
         derived in Theorem~\ref{thm:sensitivity}.}
\label{tab:data}
\begin{tabular}{@{}llrrrrrrr@{}}
\toprule
Window & Province & $\kw$ & $\nw$ & $c$ & $\mu = C\nw/N$ & $c/\mu$ & $\RR$ & $\LLR$ \\
\midrule
W\textsubscript{01} & Dong Nai  & 200 &  500 & 24 &  5.0 & 4.80 & \placeholder{?} & \placeholder{?} \\
W\textsubscript{02} & Hanoi     &  45 &  200 &  3 &  2.0 & 1.50 & \placeholder{?} & \placeholder{?} \\
W\textsubscript{03} & Ha Giang  &   3 &   30 &  2 &  0.3 & 6.67 & \placeholder{?} & \placeholder{?} \\
\midrule
National            & ---       & 2{,}800 & 10{,}000 & 100 & --- & --- & --- & --- \\
\bottomrule
\end{tabular}
\end{table}


% ── Differential Privacy ──────────────────────────────────────────────────────

\subsection{Differential Privacy}
\label{sec:background:dp}

\begin{definition}[$(\eps, \delta)$-Differential Privacy~\cite{dwork2006}]
A randomized mechanism $\mathcal{M}: \mathcal{D} \to \mathcal{R}$ is
$(\eps, \delta)$-differentially private if for all neighboring datasets
$D, D' \in \mathcal{D}$ differing in one record and all $S \subseteq \mathcal{R}$:
\[
    \Pr[\mathcal{M}(D) \in S] \leq e^\eps \cdot \Pr[\mathcal{M}(D') \in S] + \delta
\]
When $\delta = 0$ this is \emph{pure} $\eps$-DP.
\end{definition}

The \emph{global sensitivity} of a function $f$ measures the maximum change one
record can cause in the output:
\[
    \Delta f = \max_{D, D' : d(D,D') = 1} \|f(D) - f(D')\|_1
\]

\begin{definition}[$\rho$-zero-concentrated DP (zCDP)~\cite{bun2016}]
A mechanism $\mathcal{M}$ satisfies $\rho$-zCDP if for all neighboring $D, D'$
and all $\alpha \in (1, \infty)$:
\(D_\alpha(\mathcal{M}(D) \| \mathcal{M}(D')) \leq \rho\alpha\),
where $D_\alpha$ is the R\'{e}nyi divergence of order $\alpha$.
\end{definition}

The key composition property is that $T$ mechanisms satisfying
$\rho_1, \ldots, \rho_T$-zCDP compose to $(\rho_1 + \cdots + \rho_T)$-zCDP~\cite{bun2016},
and $\rho$-zCDP implies $(\rho + 2\sqrt{\rho \ln(1/\delta)}, \delta)$-DP for any $\delta > 0$.

% ── Exponential Mechanism ─────────────────────────────────────────────────────

\subsection{The Exponential Mechanism}
\label{sec:background:expomech}

\begin{definition}[Exponential Mechanism~\cite{mcsherry2007}]
Let $\mathcal{R}$ be a finite output set and $q: \mathcal{D} \times \mathcal{R} \to \mathbb{R}$
a utility function with global sensitivity $\Delta q$.
The exponential mechanism outputs $r \in \mathcal{R}$ with probability:
\[
    \Pr[\mathcal{M}_E(D) = r] \propto \exp\!\left(\frac{\eps \cdot q(D,r)}{2\,\Delta q}\right)
\]
This mechanism is $\eps$-differentially private.
\end{definition}

\begin{theorem}[Utility Guarantee~{\cite[Theorem~6]{mcsherry2007}}]
\label{thm:mst-utility}
Let $r_0 = \arg\max_r q(D,r)$ and $\Delta = q(D,r_0) - \max_{r \neq r_0} q(D,r)$.
Then:
\[
    \Pr[\mathcal{M}_E(D) \neq r_0] \leq |\mathcal{R}| \cdot \exp\!\left(\frac{-\eps \cdot \Delta}{2\,\Delta q}\right)
\]
\end{theorem}

The post-processing theorem~\cite[Proposition~2.1]{dwork2013} ensures that applying
any function to the output of a DP mechanism is still DP at no additional privacy cost.
In particular, computing the Monte Carlo p-value after the private window selection
costs no additional privacy budget.
